\documentclass[12pt,a4paper]{article}


\usepackage[a4paper,margin=2.5cm]{geometry}
\usepackage[utf8]{inputenc}
\usepackage[T1]{fontenc}
\usepackage{microtype}
\usepackage{setspace}
\setlength{\parindent}{0pt}
\onehalfspacing

% --- Content packages -------------------------------------------------------
\usepackage{graphicx}
\usepackage{amsmath,amssymb}
\usepackage{booktabs}
\usepackage{siunitx}
\usepackage{enumitem}
\usepackage{xcolor}
\usepackage{listings}
\lstset{
  language=Python,
  basicstyle=\footnotesize\ttfamily,
  keywordstyle=\bfseries,
  commentstyle=\itshape\color{black!55},
  stringstyle=\color{black!70},
  frame=single,
  rulecolor=\color{black!25},
  framesep=5pt,
  xleftmargin=6pt,
  columns=fullflexible,
  keepspaces=true,
  showstringspaces=false,
  breaklines=true,
  captionpos=b,
  aboveskip=1.2em,
  belowskip=1.2em,
}
\usepackage{tikz}
\usetikzlibrary{positioning,arrows.meta,calc}

% --- Citations --------------------------------------------------------------
\usepackage[backend=biber,style=numeric,sorting=none]{biblatex}
\addbibresource{references.bib}

% --- Links, loaded last so it wins on options ---------------------------------
\usepackage[colorlinks=true,linkcolor=blue!50!black,citecolor=blue!50!black,urlcolor=blue!50!black]{hyperref}
\usepackage{cleveref}

\title{A Large Language Model Advisory Layer for an Existing Digital Twin:\\
       \large A Tempeh Incubator Case Study}
\author{Mohamed Abdulkarim \\ Computer Engineering, Aarhus University\\[0.6em]
       \normalsize Supervisor: Cláudio Gomes}
\date{\today}

\begin{document}

\maketitle

\tableofcontents


\newpage

% =====================================================================================
\section{Introduction}
\label{sec:introduction}

\subsection{Digital twins and language models}
\label{sec:background}

A digital twin (DT) is a software system that mirrors a running physical system, with data
moving automatically in both directions between the two
\cite{fuller2020digitaltwin,kritzinger2018digitaltwin}. It is characterised by the services it
offers on top of that connection, such as monitoring, decision support and reconfiguration, and
it is not a control system: the physical system must keep working if the DT is switched off
\cite{gomes2025incubator}. This project applies the same principle one level higher up, adding a
layer on top of a DT that must keep working without that layer.

A large language model (LLM) is a neural network trained on large quantities of text that can
carry out a task described in ordinary language without being trained for it
\cite{brown2020fewshot}. Its output can be fluent and still wrong, stating false quantities with
confidence, a failure known as hallucination \cite{ji2023hallucination}. It can also call
external code and continue with the returned value \cite{schick2023toolformer}, which shifts its
responsibility from knowing a quantity to knowing which computation produces it.

The two are complementary: a DT holds a calibrated model of a process, while an LLM is easy to
address but has no reliable account of how any particular process behaves. Kleiman et al. build
a framework on this pairing, placing an agent between the user and an existing simulation
\cite{kleiman2025simulationagent}. Such a layer adds \emph{access}, turning the user's question
into a correctly configured run, and \emph{interpretation}, turning the result into an
explanation the user can act on, without the model needing to know the process itself. Two of
the framework's design choices carry over \cite{kleiman2025simulationagent}: the simulation is
called but never modified, and its raw output is reduced to compact summaries before the model
sees it, because a full time series exceeds what a model can usefully consume.


\subsection{The incubator digital twin}
\label{sec:incubator-dt}

This project builds on the incubator DT of Gomes et al. \cite{gomes2025incubator}: an insulated
box with a heating element, a fan and temperature sensors, and a controller that switches the
heater to hold the air at a setpoint. The box is used to ferment tempeh, a food made by growing a
mould on cooked soybeans over roughly a day, whose success depends on the temperature staying
within a fairly narrow range \cite{steinkraus1960tempeh,tempehinfo_temperature}. The DT consists
of thermal models of the box, calibrated against its own measurements and kept supplied with live
data, and offers three groups of services on top of them \cite{gomes2025incubator}:

\begin{itemize}
  \item \textbf{Visualisation and monitoring.} It presents the current and past state and
    compares what the models predict with what the sensors report, estimating unmeasured
    quantities from measured ones. A growing discrepancy signals that something has changed.
  \item \textbf{Decision support.} It simulates scenarios before they happen, known as what-if
    analysis, and searches across simulated configurations for the best trade-off. Gomes et al.
    give two examples: heating the box above its setpoint before a power supply is replaced, with
    simulation choosing how long to heat so that the batch neither overheats nor cools too far
    during the outage; and mapping how controller settings trade tight temperature control
    against heater wear.
  \item \textbf{Reconfiguration.} When the models no longer match the hardware, it recalibrates
    them against fresh data and re-tunes the controller.
\end{itemize}

None of these services is modified by this work.

\subsection{The gap this project addresses}
\label{sec:gap}

Consider an operator who wants to open the lid of a running batch for five minutes, and wants to
know beforehand whether the batch can take it. The DT can answer this in principle. It has
calibrated thermal models, a representation of the disturbance an open lid causes, and a
decision-support service meant for this kind of before-the-fact analysis; the power supply
example is structurally the same problem \cite{gomes2025incubator}. The gap is not a missing
capability but the distance between that capability and the person who needs it:

\begin{itemize}
  \item \textbf{Using the service means building the experiment.} A what-if analysis is a
    simulation someone configures, choosing the model, the parameters that represent the
    disturbance, the horizon and what to extract. An operator who knows fermentation but not
    thermal parameters has no route to it.
  \item \textbf{The service stops at numbers.} It returns a predicted temperature over time,
    while the operator needs a judgement about the batch. That requires knowing what the culture
    tolerates, for how long and at which stage, and the DT does not hold that knowledge: its
    models describe the air in the box, not the culture in it.
  \item \textbf{Every question has to be anticipated.} The published examples exist because
    someone foresaw those situations and wrote code for them, so what the DT can be asked is
    fixed at development time.
\end{itemize}

This project investigates a natural-language advisory layer on top of these services, which turns
the operator's question into a correctly configured call and the result back into advice about
the batch, without changing how the DT computes.

\subsection{The LLM advisory layer architecture}
\label{sec:architecture}

The advisory is organised as three ordered layers (\Cref{fig:architecture}). A layer may call the
one below it and never the one above, so a quantity reaches the operator only after something
below has computed it.

\begin{figure}[!ht]
  \centering
  \includegraphics[width=\textwidth]{llm_layer_arch.png}
  \caption{The advisory layer in three ordered layers.}
  \label{fig:architecture}
\end{figure}

The \textbf{language layer} recognises which advisory a question belongs to, converts what the
operator said into that advisory's parameters, and puts the result into words. Reading ``five
minutes'' as a duration is a unit conversion rather than a judgement, and it is the only
arithmetic the layer performs.

The \textbf{tool layer} carries out every calculation, as a pipeline. State access assembles the
batch's current conditions, the simulation projects them forward under the disturbance asked
about, and the risk assessment reduces the trajectory to a verdict, the numbers supporting it and
the margin to the nearest boundary between verdicts. Only that judgement crosses back to the
language layer, because turning a curve into a decision about the batch must not be left to
prose. The second advisory, a heater boost, reuses the same pipeline as an optimisation: each
candidate pre-heating strategy is scored by re-running the simulation and the risk assessment.

The \textbf{sources} differ in where their authority comes from. The live DT state makes an
answer concern this batch rather than batches in general. The DT's own models and controllers
are driven from outside and never altered, so numerical authority stays with the DT
\cite{gomes2025incubator}. Domain knowledge supplies the fermentation thresholds.


\newpage
\section{Implementation}
\label{sec:implementation}

\subsection{The sources}
\label{sec:sources}

Building the advisory starts at the bottom of \Cref{fig:architecture}, with the three sources
every later number is drawn from (\Cref{tab:sources}). They fix what the layers above can say,
and each can be trusted to a different degree.

\begin{table}[htbp]
    \centering
    \begin{tabular}{@{}lll@{}}
    \toprule
    \textbf{Source} & \textbf{Provides} & \textbf{Standing} \\
    \midrule
    Digital twin models & thermal and control behaviour & calibrated against the box \\
    Domain knowledge & tempeh and heater limits & sourced or estimated \\
    Current state & the running batch's condition & as published by the twin \\
    \bottomrule
    \end{tabular}
    \caption{Three sources, what each contributes, and where its authority comes from.}
    \label{tab:sources}
\end{table}

\textbf{The digital twin's models.} Four modules are taken from the twin: two plants, one
carrying the air temperature and the box's heat loss and one adding the heater element, the
controller's switching logic, and the wrapper that simulates a controller and a plant together
\cite{gomes2025incubator}. They are used exactly as the twin defines them. Their calibrated
parameters arrive with the current state rather than being written into this code, so a
recalibration on the twin's side reaches the advisory with nothing here edited.

What the models leave out bounds everything built on them. They predict the air in the box, not
the batch in it: with no term for the batch's thermal mass, a real batch cools far more slowly
than the predicted curve and never reaches its minimum, and with no term for the heat the culture
produces, a predicted drop is pessimistic during active growth. The parameters were also fitted
with the box empty and the lid closed, the open-lid disturbance is carried over from another
calibration of the same incubator without being checked against a recording, and room
temperature is held constant over the horizon.

\textbf{The current state.} A single document stands in for what the twin publishes about the
running batch: sensor readings, actuator settings, the filtered estimate of the box's
temperature, the controller's settings, the calibrated model parameters, the open-lid parameters,
and how far the batch has progressed. It is the only source that changes between one question
and the next.

\textbf{Domain knowledge.} One document holds what the twin does not (\Cref{tab:domain}). The
temperature values follow the tempeh literature
\cite{steinkraus1960tempeh,tempehinfo_temperature}; the bands, budgets and heater limits are
reasoned values adopted for this system and marked as estimated. None has been calibrated
against a batch that actually spoiled.

\begin{table}[htbp]
  \centering
  \begin{tabular}{@{}llc@{}}
    \toprule
    \textbf{Entry} & \textbf{Value} & \textbf{Basis} \\
    \midrule
    \multicolumn{3}{@{}l}{\textit{Fermentation phases, over a 24 hour batch}} \\
    \quad Lag            & 0 to 11 h, cold tolerance $\times 0.5$  & sourced \\
    \quad Active growth  & 11 to 21 h, $\times 1.75$               & sourced \\
    \quad Maturation     & 21 to 24 h, $\times 1.15$               & sourced \\
    \midrule
    \multicolumn{3}{@{}l}{\textit{Culture temperature bounds}} \\
    \quad Optimal range  & 30 to 37 \si{\celsius}, best near 34.5  & sourced \\
    \quad Growth limits  & 22 to 41 \si{\celsius}                  & sourced \\
    \quad Damage ceiling & 42 \si{\celsius}                        & sourced \\
    \midrule
    \multicolumn{3}{@{}l}{\textit{Safe air band, by phase}} \\
    \quad Lag            & 27 to 35 \si{\celsius}                  & estimated \\
    \quad Active growth  & 23 to 39 \si{\celsius}                  & estimated \\
    \quad Maturation     & 25 to 37 \si{\celsius}                  & estimated \\
    \midrule
    \multicolumn{3}{@{}l}{\textit{Time a batch may spend below its band}} \\
    \quad Still safe     & 600 s                                   & estimated \\
    \quad Recoverable    & 900 s                                   & estimated \\
    \midrule
    \multicolumn{3}{@{}l}{\textit{Heater limits}} \\
    \quad Element ceiling      & 60 \si{\celsius}                  & estimated \\
    \quad Longest run, then rest & 1800 s, then 300 s              & estimated \\
    \quad Share of any hour on & 0.75                              & estimated \\
    \bottomrule
  \end{tabular}
  \caption{Domain knowledge overview.}
  \label{tab:domain}
\end{table}



\subsection{Tool helpers}
\label{sec:helpers}

Two small modules stand between the sources and the simulation. Neither decides anything. They
make the sources available in a form the simulation can use.

\textbf{Reading the state.} The first opens the state document and validates it. It is the only
place that document is read, so a state that is malformed, incomplete or internally inconsistent
fails here with a clear message rather than reaching a simulation and producing a plausible
wrong answer.

\textbf{Reading the domain knowledge.} The second loads the thresholds in \Cref{tab:domain} and
answers two questions: which phase of fermentation the batch has reached, derived from how long
it has been running since an incubator has no sensor for it, and which temperature band counts
as safe in that phase. Both are needed before a trajectory means anything, which is why they are
settled here rather than downstream.

\subsection{The simulation tool}
\label{sec:simulation}

The simulation is the capability the language layer uses to answer a what-if question about the
lid: given an opening, it predicts what the box will do, starting from the state it is in now.

\textbf{Inputs.} The language layer supplies only the scenario: how long the lid stays open and
when it opens, by default immediately. Everything else is read from the current state, including
the starting temperatures, the room temperature, the controller settings, the calibrated
parameters and the open-lid disturbance. The language layer can therefore choose which scenario
to simulate, but not the conditions it starts from. An opening of zero length gives the undisturbed baseline,
so the same tool also answers what happens if the lid stays shut. A negative duration, an opening
longer than six hours or a state that fails validation is refused with a plain message, and
nothing is replaced by a default.

\textbf{What it does.} It runs the twin's own plant and controller together in closed loop
\cite{gomes2025incubator}, so the heater reacts to the opening as it would on the box. An open
lid is represented by multiplying the box's heat-loss coefficient, for as long as the lid is
open, by a ratio taken from the twin's more detailed calibration. The twin declares that
coefficient as an input rather than a fixed parameter, so supplying a value that changes over
time is ordinary use of its interface, and none of its equations, parameters or code is changed.
The run continues for an hour after the lid closes, long enough to observe recovery.

\textbf{What it returns.} A result has three parts. The summary holds the numbers an answer is
built from: the lowest air temperature and when it occurs, the total drop, the peak heater
temperature, the time to recover, and the time spent outside the safe band. The trajectory is
thinned to about a hundred points, which is the reduction step described in
\Cref{sec:background}; every number is still computed on the full solver grid. The inputs record
lists what was actually used, including the fermentation phase and the safe band that applied.
The tool computes and does not interpret: whether any of this is acceptable for the batch is
decided by the risk assessment (\Cref{sec:analysis}).

\textbf{Why this suits the language layer.} Three properties make the tool usable by a model
that must not compute on its own:

\begin{itemize}
  \item \textbf{It is phrased in the operator's terms.} A duration and a start time are all a call
    needs, so turning a question into a call requires no knowledge of thermal parameters. This
    removes the first obstacle in \Cref{sec:gap}.
  \item \textbf{Every number carries its assumptions.} Because the inputs record is part of the
    result, an explanation can state what a prediction rests on, such as the phase the batch was
    judged to be in, without the model reconstructing it.
  \item \textbf{It follows the twin.} Parameters and controller settings arrive with the state, so
    after the twin recalibrates or re-tunes its controller (\Cref{sec:incubator-dt}), the next
    prediction uses the new values with nothing here changed. And because a scenario is only a
    duration and a start time, a question can be explored by comparison, such as opening now
    against opening later, the same pattern of repeated calls the heater-boost advisory builds
    on (\Cref{sec:architecture}).
\end{itemize}

\newpage
\subsection{The risk assessment}
\label{sec:analysis}

The risk assessment turns the simulation's numbers into a verdict about the batch, using rules in
ordinary code so that the thresholds behind a verdict are stated rather than inferred. From the
batch's elapsed time it takes the phase, and from \Cref{tab:domain} that phase's safe air band,
its cold tolerance, and the two exposure budgets that the tolerance scales, so a young batch is
treated as less forgiving than one generating its own heat.

Time below the band within the first budget is safe, between the two budgets recoverable, and
beyond the second high risk; time above the band but under the damage ceiling of
42 \si{\celsius} turns an otherwise safe verdict into recoverable. Air above the ceiling, which
an open lid cannot cause but the heater boost could, and a run that never returns to the
setpoint are both high risk, the second because recovery cannot be confirmed, not because harm
was shown.

A recoverable verdict, the only one where what the operator does next changes the outcome, comes
with the time the lid must then stay shut. Every verdict carries the number that decided it and
the limits it inherits from the models and thresholds behind it (\Cref{sec:sources}).

For example, a five-minute opening in active growth, with the air at 30.98 \si{\celsius} and
the room at 22.4 \si{\celsius}, is predicted to take the air down to 23.81 \si{\celsius}. That
stays inside the band, so the verdict is safe, but with a margin of only 0.8 \si{\celsius}, on a
temperature the tempeh itself never reaches.


\subsection{The lid-opening advisory}
\label{sec:skill}

The language layer of \Cref{fig:architecture} holds one advisory, and it is written as
instructions rather than as code: a short document, a \emph{skill} in the terms of the assistant
that runs it, loaded when a question is about opening the lid. It names the single tool the
advisory may call and fixes the order in which the work is done, so no orchestration program of
our own sits between the operator and the tool. The assistant used here is Claude Code.

What the document fixes is how far the model's discretion reaches. The model may recognise that
a question concerns the lid, read the opening as a duration, and choose the words; it may not
fill in what the operator left out, and a question the tool will not accept returns to the
operator instead of being completed by the model. The rest is relayed: every number in an answer
must appear in the tool's output, and the caveats may be selected from but never softened. The
order is fixed as well, with the standing of the thresholds (\Cref{sec:sources}) placed before
the verdict rather than after it, so the operator cannot take a conclusion without what it rests
on.

\newpage
\printbibliography

\end{document}
